% !TeX program = pdflatex
%
% trademark-application.tex
%
% Build:  pdflatex trademark-application.tex
%

\documentclass[11pt,a4paper]{article}

\usepackage[T1]{fontenc}
\usepackage[utf8]{inputenc}
\usepackage[margin=2.4cm]{geometry}
\usepackage{amsmath}
\usepackage{amssymb}
\usepackage{booktabs}
\usepackage{longtable}
\usepackage{array}
\usepackage{xcolor}
\usepackage{enumitem}
\usepackage{fancyhdr}
\usepackage{ragged2e}
\usepackage[hidelinks]{hyperref}

\definecolor{warn}{RGB}{160,60,0}
\definecolor{ok}{RGB}{20,110,60}
\definecolor{mid}{RGB}{120,95,0}
\newcommand{\bad}[1]{\textcolor{warn}{#1}}
\newcommand{\good}[1]{\textcolor{ok}{#1}}
\newcommand{\mixed}[1]{\textcolor{mid}{#1}}
\newcommand{\code}[1]{\texttt{\small #1}}
\newcommand{\name}[1]{\textsc{#1}}

\newcolumntype{L}[1]{>{\RaggedRight\arraybackslash}p{#1}}

\pagestyle{fancy}
\fancyhf{}
\lhead{\small Isocline Trademark Application}
\rhead{\small\thepage}
\renewcommand{\headrulewidth}{0.4pt}

\title{\textbf{US Trademark Application Worksheet}\\[0.3em]
\large \name{Isocline} --- Standard Character Mark, Class 9, \S1(b)}
\author{Prepared for Starla Cortes}
\date{October 2026}

\begin{document}
\maketitle

\noindent
Draft only --- not legal advice. Have a trademark attorney review before
filing; the project's own market-positioning paper requires a clearance search
in class~9 before anything is printed or filed.

\section{Scope of this application}

\begin{center}
\begin{tabular}{@{}ll@{}}
\toprule
\textbf{Mark}     & \code{ISOCLINE} (standard character / word mark, no logo or stylization) \\
\textbf{Classes}  & 9 only (downloadable software) \\
\textbf{Basis}    & \S1(b) --- bona fide intention to use in commerce (not yet used) \\
\textbf{Filing}   & TEAS Plus (cheapest path; this application qualifies) \\
\bottomrule
\end{tabular}
\end{center}

\section{Do these before filing (in order)}

\begin{enumerate}[leftmargin=1.6em,itemsep=0.2em]
\item \textbf{Form the Ohio LLC.} The applicant must be the entity that owns
      and will use the mark. You cannot file in the name of an LLC that does
      not exist yet. Ohio SOS filing is \textasciitilde\$99, same-day online.
\item \textbf{Get the LLC's EIN} from the IRS (free, online, instant). Not
      needed on the trademark form itself, but the entity must be real and
      correctly named.
\item \textbf{Decide the LLC's exact registered name.} This exact string goes
      in the ``Owner of Mark'' field. Do not guess or abbreviate. Example
      choices and how they interact with the mark:
      \begin{itemize}[leftmargin=1.4em,itemsep=0.1em]
      \item \code{Isocline Labs LLC} --- cleanest brand-to-entity match.
      \item \code{Isocline Research LLC} --- matches the paper's research-arm
            idea, but reads more nonprofit-adjacent than a subscription app.
      \item A neutral name (\code{Cortes Software LLC}) also works; the entity
            name does not have to contain the mark.
      \end{itemize}
\item \textbf{Have a trademark attorney run a real clearance search} in
      class~9 (and ideally 42/44 as a courtesy) before filing. The Justia-only
      screen in \code{market-positioning-paper} found nothing in class~9, but
      it is explicitly \emph{not} clearance.
\end{enumerate}

\section{Mark information}

\begin{center}
\begin{tabular}{@{}L{4.4cm}L{11.6cm}@{}}
\toprule
\textbf{Field} & \textbf{Entry} \\
\midrule
Mark text & \code{ISOCLINE} \\
Mark drawing type & Standard character mark \\
Design/logo & None \\
Transliteration & None (mark is already Latin characters) \\
Translation / meaning & ``A curve along which a quantity (or slope) holds
constant.'' Do \textbf{not} claim a foreign meaning --- it is an English
technical term. \\
Disclaimer & None \\
\bottomrule
\end{tabular}
\end{center}

\noindent
A standard-character registration protects the name in any typeface, style, or
color. Do not file a logo version in this same application; a logo mark is a
separate filing. Keep this one clean.

\section{Owner of mark (applicant)}

\begin{center}
\begin{tabular}{@{}L{4.4cm}L{11.6cm}@{}}
\toprule
\textbf{Field} & \textbf{Entry} \\
\midrule
Name & \code{Isocline Labs LLC} (exact Ohio registered name) \\
Entity type & Limited Liability Company \\
State of organization & Ohio \\
Country of organization & United States \\
Citizenship statement & ``The applicant is a limited liability company
organized under the laws of Ohio.'' \\
\bottomrule
\end{tabular}
\end{center}

\section{Goods and services --- Class 9}

TEAS Plus requires the identification to track the USPTO ID Manual. Use this
wording --- functional, non-medical, and on the correct side of the
device-vs-model claim boundary:

\begin{quote}\small
\textbf{Class 9:} Downloadable computer application software for mobile devices
and tablet computers, namely, software for tracking hormone replacement therapy
injections and dosing schedules, generating medication reminders, and modeling
predicted hormone levels over time.
\end{quote}

\noindent
Alternate, more conservative ID (minimizes examiner pushback on ``hormone
replacement therapy''):

\begin{quote}\small
\textbf{Class 9:} Downloadable software in the nature of a mobile application
for tracking medication injections, scheduling medication reminders, and
generating predictions of medication levels over time.
\end{quote}

\noindent
What this application does \textbf{not} cover (do not add here):
\begin{itemize}[leftmargin=1.4em,itemsep=0.1em]
\item \textbf{Class 42} --- hosted/cloud software services (SaaS), data
      analytics. Add only if you later ship a web service or data offering.
\item \textbf{Class 44} --- medical/health services. A self-tracking app is
      \emph{not} a medical service; adding class~44 now would be rejected or
      invite an ``inaccurate identification'' objection.
\end{itemize}

\section{Filing basis --- \S1(b)}

\begin{itemize}[leftmargin=1.4em,itemsep=0.1em]
\item \textbf{Basis:} bona fide intention to use the mark in commerce on or in
      connection with the identified goods.
\item \textbf{No specimen required now.} A specimen (App Store listing
      screenshot showing the name plus a download/purchase flow) is filed with
      the Statement of Use.
\item \textbf{Deadlines:} after the Notice of Allowance you have \textbf{6
      months} to file a Statement of Use or an extension request (\S1(d)).
      Extensions run in 6-month increments up to \textasciitilde3 years (five
      extension requests). Each extension carries a filing fee.
\end{itemize}

\section{Correspondence and attorney}

\begin{center}
\begin{tabular}{@{}L{4.4cm}L{11.6cm}@{}}
\toprule
\textbf{Field} & \textbf{Entry} \\
\midrule
Attorney & Leave blank if filing \emph{pro se} (yourself) \\
Correspondence email & The LLC's email --- all USPTO mail goes here \\
Authorized to communicate & Yes (check the box permitting USPTO to email) \\
\bottomrule
\end{tabular}
\end{center}

\section{Fees (verify current amounts at uspto.gov)}

\begin{center}
\begin{tabular}{@{}ll@{}}
\toprule
\textbf{Item} & \textbf{Approximate} \\
\midrule
TEAS Plus application, one class & \textasciitilde\$350/class \\
TEAS Standard application, one class & \textasciitilde\$450/class \\
Statement of Use / extension, per class & \textasciitilde\$125/class \\
\bottomrule
\end{tabular}
\end{center}

\noindent
TEAS Plus is the right choice: standard character mark, ID-Manual-compatible
goods, no specimen yet --- all Plus requirements are met.

\section{Post-filing timeline}

\begin{enumerate}[leftmargin=1.6em,itemsep=0.1em]
\item \textbf{2--4 weeks:} serial number assigned; formality review.
\item \textbf{8--12 months:} substantive examination; expect a possible Office
      Action asking to refine the goods description.
\item \textbf{Publication:} mark published for opposition (\textasciitilde30 days).
\item \textbf{Notice of Allowance:} the 6-month statement-of-use clock starts.
\item \textbf{Statement of Use + specimen} once the app is live under the name.
\item \textbf{Registration} (\textasciitilde1 year after filing, give or take).
\end{enumerate}

\section{Open decisions}

\begin{itemize}[leftmargin=1.4em,itemsep=0.2em]
\item[$\square$] Exact Ohio LLC registered name (goes verbatim into
      ``Owner of Mark'') --- resolved: \code{Isocline Labs LLC}.
\item[$\square$] Goods description: specific (``hormone replacement therapy'')
      vs.\ general (``medication injections'') wording.
\item[$\square$] Attorney review of the class~9 clearance before filing.
\end{itemize}

\end{document}
